\chapter{Blood in Semen (Hematospermia)}

\section*{Definition}
Presence of blood in the ejaculate, presenting as pink, red, or brown-tinged semen, most often benign and self-limiting but requiring careful evaluation.

\section*{When to See a Doctor}

\begin{itemize}
 \item Arrange assessment for recurrent or persistent blood in semen, blood in men aged 40 or older, or blood with fever, urinary symptoms, pain, weight loss, or blood in the urine.
\end{itemize}

\section*{How It Develops}
Blood in semen can arise from the prostate, seminal vesicles, urethra, or other reproductive-tract structures. Infection, inflammation, recent instrumentation, and minor trauma are common explanations. Most isolated episodes are self-limited, but persistent, recurrent, or symptomatic bleeding needs more detailed assessment, especially with increasing age.

\section*{Causes}
\begin{itemize}
 \item Seminal vesicle or prostate infection (prostatitis, seminal vesiculitis)
 \item Recent prostate biopsy, instrumentation, or other genitourinary procedure
 \item Urethral stricture or trauma
 \item Prostate calculi
 \item Sexually transmitted infections
 \item Prostate or seminal vesicle cysts
 \item Hypertension or coagulopathy
 \item Prostate or other genitourinary cancer (uncommon)
 \item Testicular or prostate trauma
 \item Schistosomiasis (in endemic areas)
\end{itemize}

\section*{Related Symptoms}
\begin{itemize}
 \item pain when urinating
 \item Urinary frequency
 \item Perineal pain
 \item Fever
 \item Lower back pain
\end{itemize}

\section*{Medical Evaluation}
\begin{itemize}
 \item Your doctor will take a history including age, duration, frequency of episodes, recent genitourinary procedures, and associated urinary symptoms.
 \item Examination may include the abdomen, genitalia, testes, and prostate, guided by symptoms and age.
 \item Urinalysis and testing for sexually transmitted infections are considered when indicated. PSA testing is discussed according to age, symptoms, risk factors, and informed choice rather than used automatically.
 \item Urologic imaging, such as ultrasound or MRI, is reserved for persistent, recurrent, high-volume, or symptomatic cases or an abnormal examination.
\end{itemize}

\section*{Treatment Options}
\begin{itemize}
 \item Observation and reassurance are the mainstays for idiopathic haematospermia — spontaneous resolution occurs within weeks in most cases.
 \item Antibiotics are used only when a bacterial infection is diagnosed or strongly suspected, according to local guidance; they are not routine for isolated hematospermia.
 \item Transurethral resection is performed for obstructing prostatic cysts or calculi identified on imaging.
 \item Cystoscopy or prostate biopsy is considered only when symptoms, examination, PSA, urine findings, or imaging indicate a specific concern; age alone does not mandate these procedures.
\end{itemize}

\section*{Self-Care and Home Management}
\begin{itemize}
 \item A single isolated episode is often benign and self-limiting, but arrange assessment if it recurs or other symptoms are present.
 \item Note the colour (pink, red, dark brown) and duration of the bloody ejaculate to report to your doctor.
 \item Avoid activities that clearly trigger pain or bleeding until evaluated; there is no established need for a fixed period of abstinence in every case.
 \item Avoid vigorous sexual activity or trauma to the perineum and genitals until evaluated by a healthcare provider.
\end{itemize}

\section*{References}
\begin{thebibliography}{99}
\bibitem{blood_in_semen:ref1} Akhter W, Khan F, et al. ``Haematospermia: a systematic review.'' \textit{BJU Int}. 2013;111(4 Pt B):E163-E168.
\bibitem{blood_in_semen:ref2} Leocadio DE, Stein BS. ``Hematospermia: a review of current management.'' \textit{Urology}. 2009;73(5 Suppl):S3-S7.
\bibitem{blood_in_semen:ref3} Mathers MJ, Degener S, et al. ``Hematospermia: a diagnostic algorithm.'' \textit{Urologe A}. 2013;52(9):1251-1256.
\bibitem{blood_in_semen:ref4} Wein AJ, Kavoussi LR, et al. \textit{Campbell-Walsh Urology}, 12th ed. Elsevier, 2020.
\bibitem{blood_in_semen:ref5} American Family Physician. ``Evaluation and Treatment of Hematospermia.'' 2009;80(12):1421-1427. https://www.aafp.org/afp/2009/1215/p1421.
\end{thebibliography}
